\begin{claim}{A1.0}{Proved}{The structure-sheaf input in the precise form used below}
The direct images $R^qf_*\OO_X$, for $q=0,1,2$, are
$\OO$, $G=\OO\gamma\oplus\OO(-1)$, and $\OO(-1)$.
The inclusion $\OO\gamma\subset G$ is a subbundle.  At the cusp,
$(\gamma,u)$ is a holomorphic frame, and cup product identifies
$\bigwedge^2G$ with $R^2f_*\OO_X$.  In the cyclic covering coordinate
$s$, the action on $G$ has characters $(1,\zeta^{-a})$, where
$(m,a)=(3,2)$ or $(4,1)$ and $\zeta=e^{-2\pi i/m}$.
These assertions concern full scheme fibres, including $3S_1$ and $4S_2$.
\end{claim}
\begin{proof}
We derive the assertion from the local charts and the discriminant;
the same direct images are stated in Proposition 9.13 of
\citet{AC2026}, but that proposition is not an input to this proof.
Flatness holds
because a nonzero local parameter on the smooth base is a nonzero divisor
in the local ring of the smooth total space.  The normalization complex
of the cusp has rational surface and curve strata and the cellular
cochain complex of a torus as its constant part.  It gives
$h^q(\OO_W)=(1,2,1)$; the two degree-one classes $\gamma,u$ have nonzero
cup product.  To compute the determinant line without assuming a
canonical-bundle formula, we first retain the two finite stabilizers.
Let $\mathcal B$ be the root curve of orders three and four over $B$,
with root lines $R_1,R_2$, and put $R=R_1R_2^{-1}$.  The presentation
$3R_1=4R_2=H$ gives $\operatorname{Pic}(\mathcal B)=\Z R$,
$R_1=R^4$, $R_2=R^3$ and $\OO_{\mathcal B}(p_0)=R^{12}$;
one may also verify this by removing a line bundle's two stabilizer
characters and descending it to $\Pone$.

Write $\bar f:X\to\mathcal B$.  Translation-invariant vertical
fields on the torus charts, extended as logarithmic fields on the cusp
charts, have transition matrices depending only on the base.  They
give a rank-two bundle $V$ on $\mathcal B$ whose pullback is the
relative logarithmic tangent bundle.  Its fixed field $\xi$ is a
nowhere-zero subline, including in the logarithmic frame at the cusp.
The quotient line $Q=V/\OO\xi$ has automorphy factors $-1/\tau$
and $1/\tau$ for the two finite generators, read directly from the
linear part of the given affine deck transformations.  Thus
$\Delta(\tau)^{-1}$ is a meromorphic section of $Q^{12}$.
The discriminant product has no zero in the upper half-plane and
has a simple zero in the cusp coordinate, since
$\tau=(\log t_c)/(2\pi i)+h(t_c)$ with $h$ holomorphic.  Its divisor
therefore proves $Q^{12}=\OO_{\mathcal B}(-p_0)=R^{-12}$.
The displayed torsion-free Picard group gives
$\det V=Q=R^{-1}$.

The exact logarithmic tangent sequence has quotient
$\bar f^*T_{\mathcal B}(-p_0)$.  Its determinant, together with
$W=\bar f^*p_0$, gives
$\omega_{X/\mathcal B}=\bar f^*(\det V)^{-1}$.
The geometric fibres on the root charts have cohomology dimensions
$(1,2,1)$: they are smooth tori at the two finite points, and the
normalization calculation above gives the same result at the cusp.
Analytic base change makes the three root direct images locally free.
Connectedness gives $\bar f_*\OO_X=\OO_{\mathcal B}$, and fibrewise
Gorenstein duality with its trace pairing now yields
$R^2\bar f_*\OO_X=\det V=R^{-1}$.  The invariant local exponents
of $R^{-1}$ are two at the order-three point and one at the
order-four point.  Removing these characters leaves
$R^{-1}R_1^{-2}R_2^{-1}=R^{-12}$, so its coarse direct image is
$\OO_B(-1)$.  This supplies the degree used below independently of
the asserted direct-image or canonical-bundle formulas in \citet{AC2026}.

The same base-change calculation identifies $\gamma,u$ with a
frame through the cusp.  At a multiple fibre finite pushforward and the
Reynolds projector identify cohomology with cyclic invariants.  If $b_0$
and $b_1$ have characters $1,\zeta^{-a}$, their invariant frames are
$b_0,s^ab_1$ and $s^a(b_0\wedge b_1)$.  Their products have exactly
these exponents, so the determinant comparison is an isomorphism,
including after reduction modulo $s^m$.  Thus $\deg G=-1$.
The invariant flat class $\gamma$ does not enter the holomorphic
two-plane spanned by
$\sigma_1=6\mu\gamma+\tau u+w$ and
$\sigma_2=\beta\gamma+\mu u+\delta$: comparison of the $w,\delta$
coefficients proves this even on the smooth cyclic covers at their
centres.  It therefore gives the stated subbundle everywhere.
The quotient is $\OO(-1)$ and the extension splits because
$H^1(\OO(1))=0$.  This also specifies which subline in the otherwise
noncanonical splitting is used below.  Coherent Leray now gives
$(h^q(\OO_X))_{q=0}^3=(1,1,0,0)$, since
$H^0(\OO(-1))=H^1(\OO(-1))=H^1(\OO)=0$.
\end{proof}
